\chapter{Radial--orientation Pfaffian Hessian: diagonal no-go}

После supercritical Pfaffian gate естественно разрешить connector modulus и
проверить, может ли radial mode одновременно регулировать mass hierarchy и
orientation. Этот шаг требует нового точного расчёта: формула с \(z^{-3}\)
была справедлива только на \(|z|=1\) и не может использоваться вне circle.

\section{Self-adjoint radial continuation}

Положим
\begin{equation}
 z=re^{i\theta},
 \qquad r>0,
\end{equation}
и в обратном Dirac edge используем \(\bar z\), как требует
self-adjointness. Прямой determinant chiral block для \(W=-1/2\) равен
\begin{equation}
 \det M_-(r,\theta)
 =\frac12r^2(2r+e^{i\theta})e^{-3i\theta},
\end{equation}
поэтому
\begin{equation}
 |\det M_-|^2
 =r^4\left(r^2+r\cos\theta+\frac14\right).
\end{equation}
При \(r\to0\) determinant обращается в нуль, а Pfaffian action растёт, а не
падает. Тем самым ложная radial неограниченность исчезает.

\section{Bosonic parent potential}

Для того же полного \(J\)-парного rectangle exact traces равны
\begin{align}
 \Tr D_F^2&=12r^2+28,\\
 \Tr D_F^4&=12r^4+32r^2-8r\cos\theta+60.
\end{align}
В potential
\begin{equation}
 V_B=-\mu^2\Tr D_F^2+\lambda\Tr D_F^4
\end{equation}
условие classical stationary point \((r,\theta)=(1,0)\) фиксирует
\begin{equation}
 \frac{\mu^2}{\lambda}=\frac{13}{3}.
\end{equation}

В координатах \(\sigma=\log r\), \(\theta\) bosonic Hessian при unit vacuum
равен
\begin{equation}
 \operatorname{Hess} V_B
 =\lambda
 \begin{pmatrix}
 104&0\\0&8
 \end{pmatrix}.
\end{equation}
Reduced Pfaffian добавляет
\begin{equation}
 \operatorname{Hess}\Gamma_{\rm red}
 =
 \begin{pmatrix}
 -2/9&0\\0&2/9
 \end{pmatrix},
\end{equation}
а full KO6 modulus удваивает эту matrix.

\begin{theorem}[Diagonal radial--orientation Hessian]
Все mixed entries \(\partial_\sigma\partial_\theta\) равны нулю. Это следует
также из CP-even symmetry \(V(r,\theta)=V(r,-\theta)\). Следовательно, radial
mode не может на quadratic level динамически настроить orientation stiffness.
\end{theorem}

\section{Critical quartic coefficient}

С учётом modular orientation susceptibility \(\chi_K=0.675695\) angular
curvature около unit vacuum равна
\begin{align}
 m_{\theta,\rm red}^2&=8\lambda+\frac29-\chi_K,\\
 m_{\theta,\rm full}^2&=8\lambda+\frac49-\chi_K.
\end{align}
Spontaneous CP branch требует соответственно
\begin{align}
 \lambda&<0.056684 &&\text{(reduced)},\\
 \lambda&<0.028906 &&\text{(full KO6)}.
\end{align}
При canonical normalized value \(\lambda=1\) обе curvatures строго
положительны.

Exact two-dimensional scan полного bosonic, Pfaffian и modular functional
при \(\lambda=1\) даёт
\begin{align}
 (r_\star,\theta_\star)_{\rm red}&=(1.02426,0),\\
 (r_\star,\theta_\star)_{\rm full}&=(1.04620,0).
\end{align}
Новая flavour branch поэтому не возникает и mass/CKM test не открывается.

\begin{theorem}[Radial Pfaffian verdict]
Точный radial mode стабилен и слегка сдвигает connector modulus, но не
создаёт radial--orientation mixing. При естественной quartic normalization
vacuum остаётся CP-even. Двухмерная Pfaffian geometry сама по себе не решает
flavour trade-off.
\end{theorem}

\begin{nogocriterion}[Запрет fitted small quartic]
Нельзя уменьшать \(\lambda\) ниже critical values по CKM residuals. Ветка
переоткрывается только если spectral moments независимо выводят малую
classical quartic normalization либо parent action содержит новый CP-odd
radial--orientation invariant.
\end{nogocriterion}

Машинный ledger сохранён в
\path{s2t_v4_radial_pfaffian_hessian_gate_results.json}.